\section{Дослідження способів зберігання криптографічних ключів}

\subsection{Які загрози виникають при зберіганні криптографічних 
ключів безпосередньо у програмному коді, конфігураційних файлах або файловій системі?}

Зберігання криптографічних ключів безпосередньо у програмному коді 
або звичайних файлах створює додаткові ризики для безпеки системи. 
Наприклад, ключ може випадково потрапити до відкритого репозиторію 
GitHub разом із вихідним кодом програми або стати доступним через 
неправильні налаштування прав доступу до файлів. Також небезпеку 
становлять резервні копії та журнали роботи програм, у яких може 
залишатися конфіденційна інформація.

У разі отримання ключа зловмисник може розшифрувати захищені дані, 
підробити цифровий підпис або отримати доступ до інших ресурсів системи. 
Крім того, якщо один ключ використовується у декількох компонентах програми, 
його компрометація може вплинути на всю систему. Тому ключі не рекомендується 
зберігати у відкритому вигляді в коді або конфігураційних файлах. Для цього краще 
використовувати спеціалізовані засоби керування секретною інформацією \cite{owaspkey,owaspstorage}.

\subsection{Які можливості для захищеного зберігання ключів та іншої секретної інформації надають сучасні операційні системи та програмні платформи?}

Сучасні операційні системи мають спеціальні механізми для зберігання паролів, 
криптографічних ключів та інших конфіденційних даних. Вони дозволяють 
обмежувати доступ до секретної інформації, використовувати шифрування та визначати, 
які програми можуть працювати з певними ключами.

Наприклад, у Windows використовуються Windows Credential Manager та DPAPI, 
у macOS — Keychain, а в Android — Android Keystore. Останній, зокрема, дозволяє створювати ключі, 
які не можуть бути безпосередньо отримані програмою. На пристроях із відповідною апаратною підтримкою 
їх використання може додатково захищатися спеціальними апаратними компонентами. 
Таким чином, операційна система не просто зберігає ключі, а й контролює умови їх використання, 
що зменшує ризик викрадення секретної інформації \cite{microsoftcred,applekeychain,androidkeystore}.

\subsection{Що таке TPM (Trusted Platform Module) та для яких задач захисту криптографічних ключів він може використовуватися?}

TPM (Trusted Platform Module) — це спеціальний апаратний криптографічний модуль, призначений 
для підвищення безпеки комп'ютера. Він може бути реалізований у вигляді окремого мікрочипа 
або інтегрованого компонента. TPM дозволяє генерувати та зберігати криптографічні ключі, 
а також обмежувати умови їх використання. При цьому приватний ключ може залишатися всередині модуля, 
не передаючись програмам у відкритому вигляді.

Одним із прикладів використання TPM є захист ключів шифрування диска за допомогою BitLocker у Windows. 
Також він застосовується для автентифікації пристроїв і перевірки цілісності процесу завантаження операційної системи. 
Наприклад, використання певного ключа можна дозволити лише тоді, коли система завантажилася у визначеному довіреному стані. 
Це ускладнює отримання доступу до зашифрованих даних у разі викрадення комп'ютера \cite{microsofttpm}.

\subsection{Що таке Secure Enclave та в чому полягає основна ідея апаратно ізольованого зберігання і використання ключів?}

Secure Enclave — це спеціальне апаратно ізольоване середовище, яке використовується в пристроях 
Apple для захисту конфіденційної інформації та криптографічних ключів. Воно має окремий захищений процесор і працює 
незалежно від основної операційної системи.

Основна ідея полягає в тому, що приватний ключ створюється та зберігається всередині захищеного середовища 
і не передається основній системі у відкритому вигляді. Якщо програмі потрібно створити цифровий підпис 
або виконати іншу криптографічну операцію, вона надсилає відповідний запит до Secure Enclave, 
який самостійно виконує необхідні дії та повертає результат. Це зменшує ризик викрадення ключа 
навіть у випадку компрометації основної операційної системи. Такий механізм використовується, зокрема, 
для захисту ключів, пов'язаних з автентифікацією користувачів та цифровими підписами \cite{applesecure}.

\subsection{Що таке Hardware Security Module (HSM) та які задачі він вирішує?}

Hardware Security Module (HSM) — це спеціалізований апаратний пристрій, призначений для створення, 
зберігання та використання криптографічних ключів. Його основна особливість полягає в тому, 
що криптографічні операції виконуються безпосередньо всередині захищеного модуля, тому приватні ключі не 
потрібно передавати до звичайної пам'яті комп'ютера.

HSM використовується для шифрування та розшифрування інформації, створення цифрових підписів, 
захисту ключів центрів сертифікації та обробки платіжних операцій. Такі пристрої особливо актуальні 
для банків, великих компаній та інформаційних систем, у яких обробляється значна кількість конфіденційних даних. 
Вони також можуть проходити спеціальну сертифікацію відповідно до стандарту FIPS 140-3, який визначає вимоги до безпеки криптографічних модулів \cite{awshsm,nistfips}.

\subsection{У чому полягає різниця між зберіганням криптографічного ключа у звичайному файлі, OS Keychain/Keystore, TPM, Secure Enclave та HSM? Для яких сценаріїв можуть бути доцільними ці підходи?}

Основна різниця між розглянутими способами полягає у рівні ізоляції ключів та можливостях 
контролю доступу до них. Зберігання у звичайному файлі є найпростішим варіантом, проте він забезпечує 
найменший рівень захисту, оскільки безпека ключа значною мірою залежить від налаштувань доступу та додаткового шифрування.

OS Keychain/Keystore забезпечує системне сховище, у якому доступ до ключів контролюється операційною системою. 
Такий підхід підходить для звичайних програм і мобільних застосунків. TPM та Secure Enclave використовують апаратну ізоляцію, 
завдяки якій приватні ключі можуть залишатися всередині захищеного компонента. Відмінність полягає в тому, 
що TPM переважно використовується для захисту комп'ютерних систем, а Secure Enclave є рішенням для пристроїв Apple.

HSM, на відміну від попередніх механізмів, є спеціалізованим рішенням для організацій, яким необхідно централізовано 
керувати великою кількістю ключів та виконувати криптографічні операції з високими вимогами до безпеки.

Основні відмінності між цими підходами наведено в таблиці:

\begin{table}[h]
\centering
\caption{Порівняння способів зберігання криптографічних ключів}
\begin{tabular}{|p{3cm}|p{5cm}|p{4.5cm}|}
\hline
\textbf{Спосіб} & \textbf{Особливості} & \textbf{Сфера використання} \\
\hline
Звичайний файл & Залежить від прав доступу та шифрування файлу & Навчальні проєкти, тестування \\
\hline
OS Keychain/Keystore & Системний контроль доступу та захищене зберігання & Мобільні та настільні застосунки \\
\hline
TPM & Апаратний захист ключів, прив'язка до пристрою & Захист комп'ютерів, шифрування дисків \\
\hline
Secure Enclave & Апаратна ізоляція ключів від основної системи & Пристрої Apple, автентифікація \\
\hline
HSM & Спеціалізований апаратний захист і виконання криптографічних операцій & Банки, платіжні системи, великі організації \\
\hline
\end{tabular}
\end{table}

Отже, вибір способу зберігання залежить від призначення системи та вимог до її безпеки. Для звичайних застосунків достатньо 
можливостей Keychain або Keystore, тоді як для захисту важливих системних ключів доцільно використовувати TPM, Secure Enclave або HSM. 
Водночас навіть апаратний захист не усуває всіх ризиків, тому важливо також контролювати доступ до ключів та своєчасно виконувати 
їх заміну у разі компрометації \cite{owaspkey,owaspstorage}.

\subsection{Що означають поняття exportable та non-exportable key? У яких випадках доцільно забороняти експорт закритого ключа?}

Поняття exportable key (експортований ключ) означає криптографічний ключ, який може бути
витягнутий зі сховища ключів або криптографічного модуля у відкритому чи зашифрованому
(wrapped) вигляді для перенесення в інше середовище, резервного копіювання або використання
прикладним кодом. Натомість non-exportable key (неекспортований ключ) генерується безпосередньо
всередині захищеного середовища (HSM, TPM, Secure Enclave або ізольованого сховища ОС) і не може
покинути його межі через штатні інтерфейси ані у відкритому, ані в зашифрованому вигляді, оскільки всі
криптографічні операції з ним виконуються виключно всередині цього модуля~\cite{owaspkey, nistfips}.

Забороняти експорт закритого ключа доцільно у випадках, коли його копіювання чи компрометація
призводить до критичних або незворотних наслідків для безпеки всієї системи. Насамперед це стосується
ключів центрів сертифікації (Root CA та Intermediate CA), компрометація яких підриває довіру до всієї
інфраструктури відкритих ключів, а також ключів створення кваліфікованого або юридично значущого
електронного підпису, де критично важливо забезпечити принцип неспростовності
(non-repudiation)~\cite{nist57}. Крім того, режим non-exportable застосовують для кореневих
ключів шифрування (Master Keys) у платіжних системах, апаратних ключів автентифікації пристроїв та в
середовищах, де навіть системний адміністратор сервера не повинен мати можливості скопіювати
секрет~\cite{owaspkey}.

\subsection{Що таке PKCS\#11 та яку проблему вирішує цей стандарт?}

\subsection{Як програмний застосунок може виконати операцію цифрового підпису за допомогою HSM, не отримуючи значення закритого ключа?}

\subsection{Для чого використовується ротація криптографічних ключів? Які проблеми необхідно враховувати під час заміни ключа, яким уже зашифровано дані?}

\subsection{У яких випадках необхідне резервне копіювання ключової інформації та які додаткові ризики при цьому виникають?}

\subsection{Які стандарти регулюють вимоги до безпеки реалізації криптографічних модулів?}
